
\documentclass[11pt]{article}
\usepackage[margin=1in]{geometry}
\usepackage{amsmath,amssymb}
\usepackage{enumitem}
\usepackage[T1]{fontenc}
\usepackage{lmodern}
\setlength{\parskip}{0.5em}
\setlength{\parindent}{0pt}

\begin{document}

\begin{center}
{\LARGE \textbf{IB Physics 2016--2020}}\\
{\large \textbf{Mapped Collection for E1}}\\
\vspace{0.3em}
\textit{Faithful paraphrases of old-syllabus questions matched to the modern atomic-structure unit: Rutherford scattering, Bohr model, spectra, photoelectric effect, and wave-particle ideas.}
\end{center}

\textbf{Mapping used.} Since 2016--2020 papers predate the new syllabus, this file maps modern E1 to the closest old-syllabus material on discrete energy and radioactivity, atomic spectra, Bohr/Schr\"odinger models, photoelectric effect, and the interaction of matter with radiation.

\section*{Problems}

\begin{enumerate}[leftmargin=*, itemsep=1.2em]

\item \textbf{M17 TZ2 HL Paper 2, Q5(d) --- emission spectrum of helium}\\
Rutherford and Royds collected helium produced from alpha decay and identified it from its emission spectrum.
\begin{enumerate}[label=(\alph*)]
\item Explain, with reference to quantized atomic energy levels, how an emission spectrum is produced.
\item State why such a spectrum can be used to identify a gas.
\end{enumerate}

\item \textbf{M18 TZ2 HL Paper 2, Q9(a) --- Rutherford atomic model}\\
Rutherford used alpha-particle scattering to construct a model of the atom.
\begin{enumerate}[label=(\alph*)]
\item Describe Rutherford's atomic model.
\item State which features of scattering data led to this model.
\end{enumerate}

\item \textbf{M18 TZ2 HL Paper 2, Q9(b) --- Bohr's modification of Rutherford's model}\\
Bohr modified the Rutherford model by imposing the condition
\[
mvr=\frac{nh}{2\pi}.
\]
\begin{enumerate}[label=(\alph*)]
\item Explain why this modification was introduced.
\item State how it resolves the instability problem in the classical orbit model.
\end{enumerate}

\item \textbf{M18 TZ2 HL Paper 2, Q9(c) --- electron speed and orbit in hydrogen}\\
For the hydrogen atom, the electron is modelled in a circular orbit of radius $r$ and speed $v$.
\begin{enumerate}[label=(\alph*)]
\item Starting from Coulomb attraction providing centripetal force, derive a relation between $v$ and $r$.
\item For the state with quantum number $n=6$, calculate the electron speed.
\item Hence determine the radius of the orbit.
\end{enumerate}

\item \textbf{M19 TZ1 HL Paper 2, Q2 --- Rutherford scattering and nuclear size}\\
A scattering experiment is used to investigate the size of a thorium nucleus.
\begin{enumerate}[label=(\alph*)]
\item Use scattering data to estimate the nuclear radius of thorium-232.
\item Explain why a beam of electrons is better than a beam of alpha particles of the same energy for probing nuclear size.
\item Explain why deviations from simple Rutherford scattering are observed when very high-energy alpha particles are incident on nuclei.
\end{enumerate}

\item \textbf{N19 HL Paper 2, Q8(a)--(c) --- Bohr model, hydrogen orbits, and Schr\"odinger model}\\
A question on hydrogen compares the Bohr picture with later quantum ideas.
\begin{enumerate}[label=(\alph*)]
\item Show, by equating centripetal force and electrostatic attraction, the relation used for an electron in a Bohr orbit.
\item Use the result for $n=3$ to calculate the electron speed and its de Broglie wavelength.
\item Determine the number of wavelengths that fit around the orbit.
\item Compare the Bohr and Schr\"odinger models, including quantized energy levels and the role of the wavefunction/probability interpretation.
\end{enumerate}

\item \textbf{N19 HL Paper 2, Q11 --- photoelectric effect}\\
A graph of maximum electron kinetic energy against incident radiation frequency or wavelength is used to test competing models of light.
\begin{enumerate}[label=(\alph*)]
\item Explain why, in a classical wave picture, low-intensity light would be expected to delay electron emission.
\item Explain how the photon model accounts for immediate emission.
\item Using graph data, determine the work function of the metal.
\item Sketch how the graph changes when the intensity is reduced but the frequency range is unchanged.
\end{enumerate}

\item \textbf{N20 HL Paper 2, Q10 --- electron diffraction and nuclear size}\\
Electrons are used to investigate nuclear structure.
\begin{enumerate}[label=(\alph*)]
\item Determine the de Broglie wavelength of the electrons from their momentum or kinetic-energy data.
\item Explain how the observed scattering graph suggests diffraction.
\item Use diffraction geometry to determine the scattering angle.
\item Explain why electrons with much longer wavelength would not be suitable for resolving nuclear size.
\item Show why nuclear density is approximately independent of nucleon number $A$.
\end{enumerate}

\end{enumerate}

\end{document}
